% Original linear-Gaussian cancellation example; all noise variances are one.
\begin{tikzpicture}[x=1mm,y=1mm,font=\small,
  pgm variable/.style={circle,draw=BookInk,fill=BookPaper,
    minimum size=8mm,inner sep=0pt},
  edge/.style={-{Stealth[length=2mm]},draw=BookInk,thick}]
  \node[pgm variable] (x) at (8,8) {$X$};
  \node[pgm variable] (z) at (45,32) {$Z$};
  \node[pgm variable] (y) at (82,8) {$Y$};
  \draw[edge,BookTeal] (x)--node[above left] {$1$}(z);
  \draw[edge,BookTeal] (z)--node[above right] {$1$}(y);
  \draw[edge,BookAmber] (x)--node[below] {$-1$}(y);
  \node[font=\footnotesize,text=BookMuted] at (45,-9)
    {$Y=-X+(X+\varepsilon_Z)+\varepsilon_Y
      =\varepsilon_Z+\varepsilon_Y$};
\end{tikzpicture}
